\clearpage
\begin{examplebox}
\textbf{例1（最典型：已知焦距比，求离心率取值范围）。}

设双曲线上存在点 $P$，满足 $PF_1=2PF_2$，求 $e$ 的范围。

\textbf{解：} $\lambda=2$，所以
\[
\boxed{1<e\le\frac{2+1}{|2-1|}=3}.
\]
\textbf{等号能取到：}令 $a=1,c=3,b=2\sqrt2$，
点 $P=(1,0)$ 在双曲线上，且 $PF_1=4$、$PF_2=2$，比为 $2$。

\medskip
\textbf{例2（焦距比小于 $1$：分支不可忽略）。}

设 $PF_1=\tfrac12PF_2$，且 $P$ 在双曲线上，求离心率范围并判断 $P$ 所在支。

\textbf{解：}$\lambda=\tfrac12$，
\[
\boxed{1<e\le\frac{1+1/2}{1-1/2}=3}.
\]
因为 $PF_1<PF_2$，点 $P$ 必在\textbf{左支}；
$e=3$ 在左顶点 $(-a,0)$ 取到。
\textbf{对称性：}将 $\lambda$ 换成 $1/\lambda$，离心率范围\textbf{不变}，只交换左右支。
\end{examplebox}

\clearpage
\begin{examplebox}
\textbf{例3（从离心率反问焦距比的全部范围）。}

双曲线离心率为 $e=2$，求所有可能的 $PF_1/PF_2$。

\textbf{解：}代入固定 $e$ 的比值范围，得到
\[
\boxed{\frac{PF_1}{PF_2}\in\left[\frac13,1\right)\cup(1,3]}.
\]
左右顶点分别取 $1/3$ 与 $3$；点走向无穷远时趋近 $1$，但等于 $1$ 的点不存在。

\medskip
\textbf{例4（存在性检验：给定曲线与比例）。}

设 $H:\dfrac{x^2}{4}-\dfrac{y^2}{12}=1$，
是否存在 $P$ 使 $PF_1=2PF_2$？若存在求全部坐标。

\textbf{解：}$a=2,b=2\sqrt3,c=4,e=2<3$，故存在\textbf{两个}点。
右支焦半径 $PF_1=2x+2$、$PF_2=2x-2$；解
\[
\frac{2x+2}{2x-2}=2\quad\Longrightarrow\quad x=3.
\]
再代入双曲线：
\[
\frac94-\frac{y^2}{12}=1\quad\Longrightarrow\quad y^2=15.
\]
所以
\[
\boxed{P=(3,\sqrt{15})\text{ 或 }(3,-\sqrt{15})}.
\]
逐一核对：$PF_1=\sqrt{7^2+15}=8$，$PF_2=\sqrt{1+15}=4$，比值为 $2$。
\end{examplebox}

\clearpage
\begin{examplebox}
\textbf{例5（存在多少个点：边界分类）。}

若 $PF_1=3PF_2$，讨论 $e=\frac32,2,3$ 三种情形的解点数。

\textbf{解：}比例给出上界
\[
\frac{\lambda+1}{\lambda-1}=\frac{3+1}{3-1}=2.
\]
故 $e=\frac32$ 时为\textbf{两个}，$e=2$ 时为\textbf{一个顶点}，
$e=3$ 时为\textbf{零个}。

\medskip
\textbf{例6（高价值变式：两个焦距比互为倒数）。}

双曲线上同时存在点 $P,Q$ 满足
$PF_1=\lambda PF_2$、$QF_1=\lambda^{-1}QF_2$（$\lambda\ne1$）。
问额外增加了什么离心率限制？

\textbf{解：}\textbf{没有额外限制}。
因双曲线关于 $y$ 轴对称，$P$ 的镜像 $Q$ 交换到两焦点的距离，
自动使比例从 $\lambda$ 变成 $1/\lambda$。
两个条件均等价于
\[
\boxed{1<e\le\frac{\lambda+1}{|1-\lambda|}}.
\]
\textbf{迁移信号：}涉及对称焦点、互为倒数的焦距比，先用关于 $y$ 轴反射，
避免重复运算。
\end{examplebox}

\clearpage
\begin{examplebox}
\textbf{例7（参数表达：用 $e$ 直接寻找满足焦距比的点）。}

已知 $e=2$、$a=2$、$\lambda=\tfrac12$，求点 $P$ 的全部坐标。

\textbf{解：}$c=ea=4$，焦距比小于 $1$，故 $P$ 在左支。
设 $x=-t$，根据常规法
\[
\frac{2t-2}{2t+2}=\frac12
\quad\Longrightarrow\quad 4t-4=2t+2
\quad\Longrightarrow\quad t=3.
\]
因此 $x=-3$，$b^2=c^2-a^2=12$，
\[
\boxed{P=(-3,\sqrt{15})\text{ 或 }(-3,-\sqrt{15})}.
\]
核对：$PF_1=4, PF_2=8$，比为 $1/2$。

\medskip
\textbf{例8（考场反例：把 $\lambda=1$ 误判为任何离心率都可）。}

若题设 $PF_1=PF_2$ 且 $P$ 在非退化双曲线上，问 $e$ 范围。

\textbf{解：}\textbf{不存在这样的点，也就没有任何允许的 $e$}。
因为 $|PF_1-PF_2|=0$ 与双曲线定义 $2a>0$ 矛盾。
不能把 $\lambda=1$ 代入含 $|1-\lambda|$ 的分母，更不能将上界视作无穷而忽略不可能性。
\end{examplebox}


\clearpage
\begin{examplebox}
\textbf{例9（阿波罗尼斯圆：把“固定比例”变成圆，再判交点）。}

取双曲线 $H:\dfrac{x^2}{4}-\dfrac{y^2}{12}=1$，
焦点 $F_1(-4,0),F_2(4,0)$，规定 $PF_1=2PF_2$。
求定比轨迹圆，判定它与双曲线的交点数并求交点。

\textbf{第1步：}由 $\lambda=2,c=4$，阿波罗尼斯圆为
\[
\boxed{\mathcal A_2:\ \left(x-\frac{20}{3}\right)^2+y^2
=\left(\frac{16}{3}\right)^2}.
\]
这\textbf{不是}以 $F_1,F_2$ 为圆心的两个圆，而是所有 $PF_1/PF_2=2$ 的点\emph{共同组成的一条圆}。

\textbf{第2步：}圆心 $C=(20/3,0)$，半径 $R=16/3$，
$|CF_2|=8/3$，于是
\[
PF_2\in[R-|CF_2|,R+|CF_2|]=\left[\frac83,8\right].
\]
圆上 $PF_1-PF_2=PF_2$，双曲线需要差为 $2a=4$，
且 $4\in(8/3,8)$，故\textbf{恰有上下对称两个交点}。

\textbf{第3步：}直接解得 $x=3$，代回双曲线有 $y^2=15$，即
\[
\boxed{P_\pm=(3,\pm\sqrt{15})}.
\]

此处若改成 $a=1,c=3$（仍取 $\lambda=2$），
则 $e=3$，刚好在顶点\textbf{相切}；若 $a=1,c=4$，
则 $e=4>3$，\textbf{圆与双曲线不相交}。

\textbf{迁移价值：}给定点到两个定点的距离比时，
可先画阿波罗尼斯圆，再结合另一个\emph{轨迹方程}判断交点数。
\end{examplebox}
